\par\begingroup\small
\indent Weitere Spuren finden wir in den Tagebuchnotizen Nr. 83 vom April 1798 (eine asymptotische Formel), Nr. 86 vom Mai 1798 (die \textsc{Lagrange}sche Reihe), Nr. 87 vom Juni 1798 (die allgemeinen trigonometrischen Reihen), Nr. 88 vom 17. Juni 1798 (Beschäftigung mit \textsc{Laplace}), Nr. 104 vom 27. April 1800 (einen Satz über die Konvergenz trigonometrischer Reihen), Nr. 113 vom 25. Oktober 1800 (Hinweis auf \textls[120]{frühere} Beschäftigung mit einer asymptotischen Formel, vergl. eine Aufzeichnung in dem als Scheda A b bezeichneten Hefte des Nachlasses, die vom 5. Februar 1799 datiert ist, und den Brief an \textsc{Laplace} vom 30. Januar 1812, oben S. 371); zu allen diesen Notizen vergleiche man die Bemerkungen in dem unten folgenden Abdruck des \textit{Tagebuchs}. Endlich ist hier noch die briefliche Mitteilung zu erwähnen, die \textsc{Gauss} im Frühjahr 1799 an \textsc{v. Zach} in bezug auf die Zeitgleichungstafel von \textsc{Ulugh-Beigh}\textsuperscript{†}) hat gelangen lassen und die \textsc{Gauss} in seinen Briefen an \textsc{Olbers} vom 27. Januar 1812, Werke VIII, S. 140, und an \textsc{Schumacher} vom 3. Dezember 1831, ebenda, S. 138 erwähnt. Eine Aufzeichnung mit der Überschrift
\par\medskip\noindent\rule{.24\linewidth}{.4pt}\par\smallskip
\noindent\hspace*{\parindent}\textsuperscript{†}) Ein Auszug aus dieser Tafel war in den Allgemeinen Geographischen Ephemeriden \textbf{3}, 1799, 2. Stück, Februar S. 182, 183 in einem Berichte von \textsc{Burckhardt} über die Tafeln des \textsc{Ulugh-Beigh} erschienen; S. 181 a. a. O. heißt es: »Die 8. Tafel enthält die Mittelpunkts Gleichung; sie ist für jede 6 Minuten der Anomalie bis auf Tertien berechnet«.
\clearpage\SourceLeaf{452}\ArithmeticHead{445}{R}{BEMERKUNGEN ZUR REIHENLEHRE. BRIEFWECHSEL.}
\noindent »\textls[120]{Mittelpunktsgleichung nach \textsc{Ulugh Bey} in Zeittertien}«, in der \textsc{Gauss} die erwähnte Zeitgleichungstafel benutzt, findet sich im Nachlaß (Astr. d. 7, Kapsel 90)\textsuperscript{*}).

\indent \textsc{Jacobi} sagt (a. a. O., \textsc{Jacobis} Werke VII, S. 175), es hätte vor \textsc{Carlini} (1817) kein anderer Mathematiker oder Astronom die Aufgabe behandelt, den Koeffizienten des Sinus eines sehr großen Vielfachen der mittleren Anomalie in der Entwicklung der Mittelpunktsgleichung annähernd zu bestimmen, eine Aufgabe, die zu den schwierigsten ihrer Art gehöre. Natürlich konnte \textsc{Jacobi} nicht wissen, daß \textsc{Gauss} schon lange vor \textsc{Carlini} diese Aufgabe \textls[120]{richtig} gelöst hatte (während ja \textsc{Carlinis} Lösung fehlerhaft war) und so zu einem Ergebnis gelangt war, das mit dem übereinstimmt, das \textsc{Jacobi} selbst etwa 50 Jahre später auf dem von \textsc{Carlini} vorgezeichneten Wege von neuem abgeleitet hat. —

\par\medskip\indent In der Briefstelle [4.] spricht \textsc{Gauss} von seiner Geneigtheit, seine Methode in einer ihm »selbst genügenden Gestalt auszuführen«, und bemerkt, daß die Arbeit »einen größeren Umfang erhalten« würde; auch \textsc{Sartorius v. Waltershausen} berichtet\textsuperscript{**}), daß \textsc{Gauss} nach seinem Doktorjubiläum (16. Juli 1849) »mit der Theorie der Konvergenz der Reihen« beschäftigt gewesen sei. Wir werden also die Entwürfe [V.] und [VI.] als Vorarbeiten zu der Abhandlung anzusehen haben, deren Plan \textsc{Gauss} in der Briefstelle [4.] erwähnt. Diese Entwürfe sind also nach dem 5. Februar 1850 entstanden\textsuperscript{***}), und wir können nunmehr über das, was jene geplante Abhandlung enthalten sollte, die folgenden zusammenfassenden Angaben machen. Auf eine Theorie der komplexen Größen sollte die Entwicklung der Lehre von dem Zusammenhang der Flächen (Analysis situs) folgen, die bestimmt war, der Lehre von der Integration im komplexen Gebiete als Grundlage zu dienen. Wahrscheinlich würde sich \textsc{Gauss} aber auch hier — ähnlich wie in dem Briefe an \textsc{Bessel} vom 18. Dezember 1811 (siehe oben S. 365, besonders S. 367 Fußnote) — auf die Integrale \textls[120]{einförmiger} Funktionen beschränkt haben, wenngleich die von \textsc{Dedekind} berichtete Äußerung (oben S. 438, Fußnote) auch der gegenteiligen Möglichkeit Raum gibt. Eine Erörterung der Eigenschaften allgemeiner (auch komplexer) Zahlenfolgen würde zu dem eigentlichen Gegenstande der Abhandlung übergeleitet haben, zur Bestimmung des infinitären Verhaltens der Koeffizienten einer allgemeinen trigonometrischen Reihe und der Untersuchung ihrer Konvergenz. Die Anwendung auf die Entwicklung der Mittelpunktsgleichung hätte den Abschluß gebildet\textsuperscript{†}).

\indent Wie \textsc{Gauss} die \textls[120]{allgemeinen} trigonometrischen Reihen zu behandeln gedachte, dafür gibt der Nachlaß kaum einen Fingerzeig; daß bei der Mittelpunktsgleichung die imaginäre Substitution (siehe oben die Gleichungen (\(\ast\)) auf S. 439)
\begin{equation}
E=i\log\operatorname{cotang}\frac\varphi2+\varepsilon\tag*{(**)}
\end{equation}
\par\medskip\noindent\rule{.24\linewidth}{.4pt}\par\smallskip
\noindent\hspace*{\parindent}\textsuperscript{*}) In dem erwähnten Briefe an \textsc{Olbers} vom 27. Januar 1812, Werke VIII, S. 140, schreibt \textsc{Gauss}, daß die Papiere, worin er 1799 die Methode der kleinsten Quadrate auf \textsc{Ulugh-Beighs} Zeitgleichungstafel angewandt hatte, verloren gegangen seien; allem Anschein nach hat \textsc{Gauss} sie 1850, zugleich mit den auf die Mittelpunktsgleichung bezüglichen Aufzeichnungen, wiedergefunden.
\par\noindent\hspace*{\parindent}\textsuperscript{**}) \textit{Gauss zum Gedächtniss}, 1856, S. 69.
\par\noindent\hspace*{\parindent}\textsuperscript{***}) Im art. 2. des Abschnitts [VI.], oben S. 408, wird auf die \textit{Jubiläumsschrift} von 1849 Bezug genommen.
\par\noindent\hspace*{\parindent}\textsuperscript{†}) Nach den Akten der Philosophischen Fakultät zu Göttingen (vol. 117, S. 84) hat \textsc{Gauss} im Jahre 1833 drei Preisaufgaben vorgeschlagen, von denen die eine den Wortlaut hatte: »Enarrentur variae methodi problema \textsc{Kepleri} solvendi, imprimis per series infinitas revoceturque gradus convergentiae, quam hae offerunt, ad mensuram accuratam«. Die Fakultät hat aber eine andere der vorgeschlagenen Aufgaben (nämlich, die Trägheitsmomente der regelmäßigen Körper zu bestimmen) gewählt und gestellt.
\clearpage\SourceLeaf{453}\ArithmeticHead{446}{L}{ANALYSIS. NACHLASS.}
\noindent sofort auf Integrale im komplexen Gebiet führt, wenn man die Koeffizienten der Entwicklung in Integralform darstellt, ist einleuchtend. Die Bedeutung dieser Substitution erläutert J. \textsc{Horn} in einer brieflichen Mitteilung wie folgt:

\indent »Für die Entwicklung (III), oben S. 439, ist
\[
C_n=\frac1\pi\int_{-\pi}^{+\pi}\frac{dv}{dM}e^{inM}\,dM
=\frac{\cos\varphi}\pi\int_{-\pi}^{+\pi}e^{inM}\frac{dE}{1-f\cos E},
\]
\noindent oder, indem man\textsuperscript{*}) \(z=e^{iE}\) einführt,
\[
C_n=-\frac{i\cos\varphi}\pi\int F(z)(\Phi(z))^n\,dz,
\]
\noindent wobei
\[
F(z)=\frac1{z-\frac f2(z^2+1)},\quad
\Phi(z)=ze^{-\frac f2\left(z-\frac1z\right)}
\]
\noindent gesetzt ist und über den Kreis \(|z|=1\) integriert wird. Die Funktion \(F(z)\) hat die reellen singulären Stellen
\[
\begin{aligned}
z_0&=\frac{1-\sqrt{1-f^2}}f=\operatorname{tang}\frac\varphi2=e^{i\mathfrak E}<1,\\[4pt]
z_1&=\frac{1+\sqrt{1-f^2}}f=\operatorname{cotang}\frac\varphi2=e^{-i\mathfrak E}>1,
\end{aligned}
\]
\noindent die zugleich Nullstellen von \(\Phi'(z)\) sind. Der Integrationsweg \(|z|=1\) kann durch den Integrationsweg \(|z|=z_0\) ersetzt werden, der nur der singulären Stelle \(z_0\) ausweichen muß. Demgemäß hat man zu setzen
\[
z=z_0e^{i\varepsilon}
\]
\noindent oder, was dasselbe ist,
\[
E=\mathfrak E+\varepsilon=i\log\operatorname{cotg}\frac\varphi2+\varepsilon,
\]
\noindent und dies ist eben die von \textsc{Gauss} angewandte Substitution. Nimmt man als Integrationsweg den Kreis \(|z|=z_1\) mit Umgehung der singulären Stelle \(z_1\), so hat man
\[
z=z_1e^{\vartheta i}\ \text{oder}\ E=\vartheta-i\log\operatorname{cotg}\frac\varphi2,
\]
\noindent eine Substitution, die in einer Abhandlung von W. \textsc{Scheibner}\textsuperscript{**}) benutzt wird.«

\indent Soweit die Mitteilung von \textsc{Horn}\textsuperscript{***}).
\par\medskip\noindent\rule{.24\linewidth}{.4pt}\par\smallskip
\noindent\hspace*{\parindent}\textsuperscript{*}) Vergl. für das folgende H. \textsc{Burkhardt}, \textit{Über Funktionen grosser Zahlen, insbesondere über die näherungsweise Bestimmung entfernter Glieder in den Reihenentwicklungen der Theorie der Keplerschen Bewegung}, Sitzungsberichte der Kgl. Bayerischen Akademie der Wissensch. 1914, S. 1.
\par\noindent\hspace*{\parindent}\textsuperscript{**}) W. \textsc{Scheibner}, \textit{Über die asymptotischen Werthe der Coefficienten in den nach der mittleren Anomalie vorgenommenen Entwicklungen}, zuerst erschienen in den Berichten der Kgl. Sächsischen Gesellsch. der Wissensch. zu Leipzig, Mathem.-phys. Classe \textbf{8}, 1856, S. 40, dann gekürzt in den Mathematischen Annalen \textbf{17}, 1880, S. 545, siehe hier S. 551.
\par\noindent\hspace*{\parindent}\textsuperscript{***}) Man vergl. auch die in der erwähnten Abhandlung von \textsc{Burkhardt} angeführte Literatur, der noch hinzuzufügen wäre A. L. \textsc{Cauchy}, \textit{Mémoire sur divers points d'Analyse}, Mémoires de l'Académie des Sciences \textbf{8}, Paris 1829, S. 101, Oeuvres 1. série, II, S. 33.
\clearpage\SourceLeaf{454}\ArithmeticHead{447}{R}{BEMERKUNGEN ZUR REIHENLEHRE. BRIEFWECHSEL.}
\indent In dem auf S. 382 beginnenden Hauptstück haben wir das im Nachlaß befindliche Material gesammelt, das mit dieser unvollendet gebliebenen Abhandlung von \textsc{Gauss} im Zusammenhang steht. Der Brief [5.] gehört hierher, weil \textsc{Gauss} sich darin über seine Auffassung der Reihenkonvergenz ausspricht, der Brief [6.], weil darin von der »Metaphysik der komplexen Größen« gehandelt wird\textsuperscript{*}). Der Vollständigkeit wegen wäre noch eine Stelle in dem Briefe von \textsc{Gauss} an \textsc{Wolfgang Bolyai} vom 20. April 1848\textsuperscript{**}) zu nennen, wo von einem allgemeinen Kriterium für die Konvergenz der Reihensummen die Rede ist.
\par\hfill\textsc{Schlesinger}.
\par\medskip\noindent\rule{.24\linewidth}{.4pt}\par\smallskip
\noindent\hspace*{\parindent}\textsuperscript{*}) Die Urschrift des an \textsc{Grassmann} gerichteten Briefes befindet sich im Besitz der Familie \textsc{Grassmanns}; der Brief ist bereits abgedruckt bei V. \textsc{Schlegel}, \textit{H. Grassmann, sein Leben und seine Werke}, Leipzig 1878, S. 22—23 und in H. \textsc{Grassmanns} Gesammelten Werken I, 2, Leipzig 1896, S. 397. Der oben stehende Abdruck ist mit der Urschrift verglichen worden.
\par\noindent\hspace*{\parindent}\textsuperscript{**}) \textit{Briefwechsel zwischen C. Fr. Gauss und W. Bolyai}, herausg. von F. \textsc{Schmidt} und P. \textsc{Stäckel}, Leipzig 1899, S. 133, 5); vergl. \textsc{Bolyais} Bemerkung, ebenda, S. 139 und die zu S. 133—134 gehörige Anmerkung \textsc{Stäckels}, ebenda, S. 195, 196.
\par\bigskip\begin{center}\rule{.24\linewidth}{.4pt}\end{center}
\endgroup
\ArithmeticHead{447}{R}{BEMERKUNGEN ZUR REIHENLEHRE. BRIEFWECHSEL.}
